\section{方法}

\subsection{统一闭环框架}

本文的方法将扩散模型放入一条任务驱动的物理闭环中，整体逆向设计流程围绕统一响应对象展开。首先把目标器件功能转写为可计算的 task-aware score，这个分数用于刻画和量化器件在目标任务上的性能表现，并据此定义训练、筛选和最终报告所依赖的目标响应对象。随后，这个统一响应对象同时驱动数据构建、模型训练、条件生成和最终基于仿真验证的重评分。若记目标响应为 $\mathbf{R}^{\star}$，则整个流程可以写成
\begin{equation}
\mathbf{R}^{\star}\xrightarrow{\ \mathcal{S}\ }\mathcal{D}_{\mathrm{train}}
\xrightarrow{\ \mathcal{F},\mathcal{G}\ }\widehat{\mathcal{G}}_{1:K}
\xrightarrow{\ \mathrm{Sim.}\ }\widehat{\mathcal{G}}_{\mathrm{best}},
\end{equation}
其中 $\mathcal{S}$ 表示任务打分函数族，它把目标 OTF 与器件响应之间的吻合程度转写为可比较的任务分数，从而直接量化器件在目标操作上的表现；$\mathcal{D}_{\mathrm{train}}$ 表示仿真驱动的数据集，$\mathcal{F}$ 表示前向代理模型，$\mathcal{G}$ 表示条件扩散生成器，$\widehat{\mathcal{G}}_{1:K}$ 是生成出的候选结构集合。最终输出由相同任务指标下得分最高、且通过仿真验证的结构确定，因为本文真正报告的是器件在目标任务上的表现，而不是网络内部的名义最优预测。

这样的组织方式有两个直接好处。第一，前向模型和逆向模型训练时面对的是同一个响应对象，最终物理验证也仍然作用在这个对象上。第二，候选集合在整个推理阶段被明确保留下来，这对一对多逆向问题至关重要。这里还需要强调前向代理的特殊作用：扩散模型本身并不直接编码超表面的电磁响应规律，它主要从数据分布中学习结构生成模式；前向代理则提供结构到响应的近似映射，并把响应一致性写成训练阶段的物理损失项。这个物理损失会把扩散采样拉回到满足目标电磁响应的区域，使生成结果同时受到数据先验和物理约束。因此，神经网络在这里承担的是仿真预算分配、目标可行集填充和物理约束注入的角色，最终裁决仍由仿真验证完成。

\subsection{数据集构建、打分和增强}

本文的几何类是 $64\times64$ 的自由形态二值图案，约定 $1=\mathrm{material}$、$0=\mathrm{air}$。结构生成器从粗网格随机场出发，施加 $C_4+\sigma_x+\sigma_y$ 对称并上采样到最终分辨率。这里采用 $D_4$ 型对称先验，是为了在保留自由形态表达能力的同时强化面内各向同性；在这一约束下，$0^{\circ}$--$45^{\circ}$ 扇区内的结构或角响应信息，基本就足以通过对称折叠确定整个超表面。随后再通过两轮高斯平滑与二值化逐步形成可制造的自由形态结构：第一轮用于将粗糙随机场平滑为连续形状原型，第二轮用于进一步规整边界并抑制细碎高频细节。为近似满足最小加工尺寸约束，代码随后引入形态学开闭运算，其中开运算去除小毛刺、小孤岛和过细桥连，闭运算填补小孔洞与窄裂缝；若清理过强，则退化为更保守的闭运算方案。这样得到的结构空间既明显超出少量解析参数族，又不会完全脱离制造友好约束。

所有结构都通过 RCWA 在固定物理堆栈下标注。这里采用的物理设置是：周期 $500$~nm、厚度 $500$~nm、结构材料为硅、入射介质为空气、输出介质为二氧化硅。响应网格包含 $11$ 个波长采样点（$800$--$1300$~nm，步长 $50$~nm）和 $17$ 个角度采样点（$-40^{\circ}$ 到 $40^{\circ}$，步长 $5^{\circ}$）。同时保留两个偏振保持通道，即 $T_{pp}$ 和 $T_{ss}$ 的幅值响应，因此每个条件输入是形状为 $[2,11,17]$ 的张量。虽然网格规模不大，但已经足以编码 wavelength 漂移、角度曲率和偏振泄漏这些算子器件最关键的响应特征。

记 $R_q(\lambda_j,\theta)$ 为偏振通道 $q\in\{p,s\}$ 在采样波长 $\lambda_j$ 处的验证响应切片，并定义按行归一化的响应为
\begin{equation}
\widetilde{R}_q(\lambda_j,\theta)=\frac{R_q(\lambda_j,\theta)}{\max_{\theta}R_q(\lambda_j,\theta)+\varepsilon}.
\end{equation}
本文所用的评分函数族，本质上是把不同物理目标投影到一组可验证、可比较的任务评价准则上。它的作用不只是仿真后排序，而是把光学规范明确转写成后续学习、增强和筛选所共同依赖的目标函数。

以本文这一兼具大NA、高透过率和宽带工作的二阶微分任务为例，其理想角度包络写成
\begin{equation}
g_2(\theta)=\left|\frac{\sin\theta}{\sin\theta_{\max}}\right|^{2},
\end{equation}
对应的行分数定义为
\begin{equation}
S_2=(1-\alpha)\left(w_c S_c+w_s S_s+w_e S_e\right)+\alpha S_b.
\end{equation}
这里 $S_c$ 用于奖励 $\theta=0$ 附近的中心抑制，$S_s$ 衡量响应与目标 $k_x^2$ 型包络的吻合程度，$S_e$ 奖励大角度处足够的透过支撑，$S_b$ 则要求相同的算子行为在邻近波长上继续保持。从物理上说，这四项分别对应低空间频率压制、正确的微分型角谱增长、高空间频率信息保留以及有限带宽稳健性。其他任务的具体评分定义遵循同样的响应中心思想，详见补充材料。

\subsection{前向代理模型}

前向代理学习的是
\begin{equation}
\mathcal{F}_{\phi}: \mathbf{G}\in\{0,1\}^{64\times64}\mapsto \widehat{\mathbf{R}}\in\mathbb{R}^{2\times11\times17},
\end{equation}
其中 $\widehat{\mathbf{R}}(:,j,k)$ 表示第 $j$ 个波长采样点、第 $k$ 个角度采样点处的双偏振联合预测。它的作用边界被刻意限定得很清楚，但物理上又十分关键：一方面它近似结构到响应场的映射，用于大规模候选预筛；另一方面它为扩散训练提供响应一致性约束，并把有限的 RCWA 预算集中到更可能满足目标响应的候选上。

从架构上看，当前 surrogate 采用的是 coordinate-aware encoder 加上 multiscale response-grid fusion。输入端在二值几何图之外显式拼接了两维笛卡尔坐标通道，使网络能够学习平面几何分布与远场角度响应之间的关系，而不是把超表面简单视为平移不变纹理。随后，多层残差卷积块逐步下采样并扩大感受野，使模型同时捕捉局部微结构散射细节，以及整体拓扑、对称性和填充分布对角谱包络与光谱漂移的影响。不同尺度特征并不重新解码回原图平面，而是统一投影到最终的 $11\times17$ wavelength--theta 标签网格上再融合，因此预测头直接工作在真实物理采样平面上。

另外两个设计点进一步补足了器件级物理信息。其一，最深层 latent 经过全局多层感知机压缩后广播回整个响应图，使每个 wavelength--theta 位置都能感知整体拓扑和总填充比等全局上下文。其二，输出侧再次注入响应网格自身的坐标信息，使网络显式区分中心角与外角、短波侧与长波侧的预测位置。两个偏振保持通道则采用联合输出而不是分别建模，因为 $T_{pp}$ 与 $T_{ss}$ 共享同一几何来源，通道相关性和泄漏本身就是重要物理信息。

训练时，对每个 channel、wavelength、theta 位置分别做标准化，从而避免强响应区域主导损失，并保持对零点、泄漏带等弱响应但物理关键位置的敏感性。同时，在统计量计算前先过滤仿真失败产生的 NaN 样本，数据增强则严格限制在保持光学响应不变的对称翻转上。代理模型的主损失是归一化响应空间中的 $L_1$ 重建：

代理模型的主损失是归一化响应空间中的 $L_1$ 重建：
\begin{equation}
\mathcal{L}_{\mathrm{fwd}}=\left\|\mathcal{F}_{\phi}(\mathbf{G})-\widetilde{\mathbf{R}}\right\|_{1}.
\end{equation}
同时，训练日志会在反归一化后追踪物理空间误差。在当前保留的训练日志中，归一化验证损失从初始约 $0.63$ 降到训练末期约 $0.366$。这一精度已经足以让前向代理成为有效的 screening oracle，尽管最终排序和结论仍然必须交由严格电磁仿真完成。

\subsection{条件扩散逆向生成}

逆向模型采用条件扩散过程：把目标响应张量作为条件输入，从中采样二值结构候选。训练时，输入结构从 $[0,1]$ 映射到 $[-1,1]$，采用 cosine noise schedule，并使用 velocity parameterization，网络学习目标由此落在速度变量上。记干净结构为 $x_0$，时刻 $t$ 的加噪样本为 $x_t$，归一化条件为 $\mathbf{c}$，则正向加噪过程为
\begin{equation}
x_t=\sqrt{\bar{\alpha}_t}\,x_0+\sqrt{1-\bar{\alpha}_t}\,\epsilon,\qquad \epsilon\sim\mathcal{N}(0,I),
\end{equation}
网络学习的是 $v_{\theta}(x_t,t,\mathbf{c})$，其基本损失为
\begin{equation}
\mathcal{L}_{\mathrm{diff}}=\left\|v_{\theta}(x_t,t,\mathbf{c})-v^{\star}(x_0,\epsilon,t)\right\|_2^2.
\end{equation}

模型主干是条件 U-Net，但当前实现里有三个点格外关键。第一，条件本身采用结构化的双通道响应场表示，因此使用了专门的 condition encoder。这个编码器会把目标响应同时映射成全局条件嵌入、分层条件特征图以及一组光谱 token，使同一个目标既能以压缩摘要形式进入网络，也能以保留网格结构的形式参与各尺度计算。第二，全局条件嵌入与时间嵌入一起注入残差块，分层条件特征图则在对应分辨率上与主干特征融合；与此同时，cross-attention 会把光谱 token 注入潜变量层级，使局部几何去噪过程始终受到整体 wavelength--theta 响应模式的约束。第三，训练时通过 condition dropout 支持 classifier-free guidance，推理时可以用 guidance 强化目标约束。

对当前问题而言，单纯的扩散损失还不够，因为视觉上像二值结构的样本未必物理上可用。因此总损失中还加入了两个辅助项：
\begin{equation}
\mathcal{L}_{\mathrm{total}}
=\mathcal{L}_{\mathrm{diff}}
+\lambda_{\mathrm{phys}}\mathcal{L}_{\mathrm{phys}}
+\lambda_{\mathrm{bin}}\mathcal{L}_{\mathrm{bin}},
\end{equation}
其中 $\mathcal{L}_{\mathrm{phys}}=\|\mathcal{F}_{\phi}(\widehat{x}_0)-\mathbf{c}\|_1$ 用于鼓励去噪结果在前向代理下保持目标响应，$\mathcal{L}_{\mathrm{bin}}$ 通过 $x(1-x)$ 惩罚灰区，使采样结果更接近真正二值。这里的 $\mathcal{L}_{\mathrm{phys}}$ 就是 2.1 所说的物理约束入口：前向代理把结构对应的响应近似映射回目标空间，并将响应一致性显式插入扩散训练。在当前保留的训练设置中，$\lambda_{\mathrm{phys}}=0.5$、$\lambda_{\mathrm{bin}}=0.05$。保存下来的验证日志显示，最终总验证损失约为 $0.129$，其中物理项始终不是可忽略的小修正。这说明扩散模型并非只在学习“像训练集”，同时也被显式拉回物理目标空间。

\subsection{物理引导推理与仿真重评分}

推理阶段采用明确的候选集策略。对于每个目标 case，先构造定义在 $(\lambda,\theta,p)$ 上的理想目标光谱 OTF，并用保存下来的统计量做标准化，然后重复成 batch 作为条件输入。扩散模型一次生成多个二值结构候选，再由 surrogate 提供快速初筛，最后把更有希望的候选送回仿真验证流程，按照任务指标重新打分。

最终指标总是任务相关的，但始终作用在验证响应上。例如二阶微分分数由中心抑制、目标偶对称曲线拟合和外角度支持共同决定；偏振无关分数加入通道匹配；偏振复用分数要求主动通道保持响应、被动通道尽量静默；时空微分分数则直接作用在二维 wavelength--theta 窗口上。这种“神经网络负责生成和预筛，仿真验证负责最后裁判”的分工，就是本文所谓物理一致性的实际含义。
